\documentclass[12pt,a4paper]{article}

\usepackage[utf8]{inputenc}
\usepackage[T1]{fontenc}
\usepackage[french]{babel}
\usepackage{geometry}
\geometry{margin=2.5cm}
\usepackage{graphicx}
\usepackage{hyperref}
\usepackage{xcolor}
\usepackage{listings}
\usepackage{booktabs}
\usepackage{fancyhdr}
\usepackage{titlesec}
\usepackage{enumitem}
\usepackage{amsmath}
\usepackage{tcolorbox}
\usepackage{mdframed}
\usepackage{float}
\usepackage{microtype}
\usepackage{lmodern}

% Colors
\definecolor{primary}{RGB}{37,99,235}
\definecolor{secondary}{RGB}{124,58,237}
\definecolor{darkgray}{RGB}{30,30,30}
\definecolor{codebg}{RGB}{245,247,250}
\definecolor{codeframe}{RGB}{200,210,230}

% Hyperref setup
\hypersetup{
    colorlinks=true,
    linkcolor=primary,
    urlcolor=secondary,
    citecolor=primary
}

% Section title formatting
\titleformat{\section}{\Large\bfseries\color{primary}}{}{0em}{\thesection\quad}[\titlerule]
\titleformat{\subsection}{\large\bfseries\color{darkgray}}{}{0em}{\thesubsection\quad}
\titleformat{\subsubsection}{\normalsize\bfseries\color{secondary}}{}{0em}{\thesubsubsection\quad}

% Code listings style
\lstset{
    backgroundcolor=\color{codebg},
    frame=single,
    rulecolor=\color{codeframe},
    basicstyle=\ttfamily\small,
    keywordstyle=\color{primary}\bfseries,
    commentstyle=\color{gray}\itshape,
    stringstyle=\color{secondary},
    numbers=left,
    numberstyle=\tiny\color{gray},
    breaklines=true,
    captionpos=b,
    tabsize=2
}

% Header/Footer
\pagestyle{fancy}
\fancyhf{}
\fancyhead[L]{\small\textcolor{primary}{\textbf{Cloud, Virtualization \& Containers}}}
\fancyhead[R]{\small\textcolor{gray}{Rapport -- 2026}}
\fancyfoot[C]{\thepage}
\renewcommand{\headrulewidth}{0.4pt}

% Tcolorbox style
\tcbuselibrary{skins,breakable}
\newtcolorbox{keybox}[1]{
    colback=blue!5!white,
    colframe=primary,
    fonttitle=\bfseries,
    title=#1
}

% -----------------------------------------------------------
\begin{document}

% ==================== TITLE PAGE ====================
\begin{titlepage}
\centering
\vspace*{2cm}
{\Huge\bfseries\color{primary} Cloud, Virtualisation \& Conteneurs\par}
\vspace{0.5cm}
{\large\color{gray}\rule{0.6\textwidth}{1pt}\par}
\vspace{1cm}
{\Large\bfseries Rapport de Module\par}
\vspace{0.5cm}
{\large Matière : Cloud, Virtualization \& Containers\par}
\vspace{2cm}

\begin{tcolorbox}[colback=blue!5,colframe=primary,width=0.75\textwidth]
\centering
\textbf{Thèmes abordés}\\[6pt]
Virtualisation Système · Conteneurisation Docker\\
Orchestration Kubernetes · Cloud Computing\\
Infrastructure as Code · CI/CD \& DevOps
\end{tcolorbox}

\vfill
{\large\textbf{Étudiant :} Yassine Daggaz\par}
{\large\textbf{Filière :} Master / Ingénierie Informatique\par}
\vspace{0.5cm}
{\large\textcolor{gray}{Année universitaire 2025 -- 2026}\par}
\end{titlepage}

% ==================== TABLE OF CONTENTS ====================
\tableofcontents
\newpage

% ==================== INTRODUCTION ====================
\section{Introduction}

Le cloud computing, la virtualisation et la conteneurisation forment aujourd'hui un triptyque fondamental de l'infrastructure informatique moderne. Ces technologies permettent d'optimiser l'utilisation des ressources matérielles, d'améliorer la portabilité des applications et d'automatiser le déploiement à grande échelle.

\begin{keybox}{Objectifs du rapport}
\begin{itemize}[noitemsep]
    \item Comprendre les principes de la virtualisation et ses différentes formes
    \item Maîtriser la conteneurisation avec Docker
    \item Explorer l'orchestration avec Kubernetes
    \item Appréhender les modèles de services Cloud (IaaS, PaaS, SaaS)
    \item Mettre en pratique ces technologies via un projet concret
\end{itemize}
\end{keybox}

% ==================== VIRTUALISATION ====================
\section{Virtualisation}

\subsection{Définition et principes}

La virtualisation est une technologie qui permet de créer une version logicielle d'une ressource physique (serveur, réseau, stockage). Elle repose sur un composant clé appelé \textbf{hyperviseur} qui gère les ressources entre les machines virtuelles (VM).

\subsection{Types de virtualisation}

\subsubsection{Virtualisation complète (Type 1 -- Bare Metal)}

L'hyperviseur s'exécute directement sur le matériel physique. Exemples : VMware ESXi, Microsoft Hyper-V, KVM.

\begin{figure}[H]
\centering
\begin{tcolorbox}[colback=gray!10,colframe=gray,width=0.7\textwidth]
\centering
\texttt{[ VM1 ] [ VM2 ] [ VM3 ]}\\
\texttt{--- Hyperviseur (Type 1) ---}\\
\texttt{--- Matériel physique ------}
\end{tcolorbox}
\caption{Architecture Hyperviseur Type 1}
\end{figure}

\subsubsection{Virtualisation hébergée (Type 2)}

L'hyperviseur s'exécute au-dessus d'un système d'exploitation hôte. Exemples : VirtualBox, VMware Workstation.

\subsubsection{Para-virtualisation}

Les VMs sont conscientes d'être virtualisées et communiquent avec l'hyperviseur via des pilotes optimisés (ex. : Xen).

\subsection{Avantages de la virtualisation}

\begin{table}[H]
\centering
\begin{tabular}{@{}lp{8cm}@{}}
\toprule
\textbf{Avantage} & \textbf{Description} \\
\midrule
Isolation & Chaque VM est complètement isolée \\
Consolidation & Réduction du nombre de serveurs physiques \\
Haute disponibilité & Migration à chaud (\textit{live migration}) \\
Snapshots & Sauvegarde instantanée de l'état \\
Multi-OS & Exécution de plusieurs OS simultanément \\
\bottomrule
\end{tabular}
\caption{Avantages de la virtualisation}
\end{table}

% ==================== CONTENEURISATION ====================
\section{Conteneurisation avec Docker}

\subsection{Conteneur vs Machine Virtuelle}

Contrairement aux VMs, les conteneurs partagent le noyau du système hôte et ne virtualisent que l'espace utilisateur. Ils sont donc beaucoup plus légers et démarrent en quelques secondes.

\begin{table}[H]
\centering
\begin{tabular}{@{}lcc@{}}
\toprule
\textbf{Critère} & \textbf{VM} & \textbf{Conteneur} \\
\midrule
Taille & Plusieurs Go & Quelques Mo \\
Démarrage & Minutes & Secondes \\
Isolation & Complète (OS) & Processus \\
Performance & Overhead élevé & Quasi-native \\
Portabilité & Limitée & Très élevée \\
\bottomrule
\end{tabular}
\caption{Comparaison VM vs Conteneur}
\end{table}

\subsection{Architecture Docker}

Docker repose sur une architecture client-serveur :

\begin{itemize}
    \item \textbf{Docker Client} : Interface CLI (\texttt{docker build}, \texttt{docker run}, ...)
    \item \textbf{Docker Daemon} : Service en arrière-plan (\texttt{dockerd})
    \item \textbf{Registry} : Dépôt d'images (Docker Hub, GitHub Container Registry)
    \item \textbf{Images} : Templates immuables en couches (\textit{layers})
    \item \textbf{Conteneurs} : Instances en cours d'exécution d'une image
\end{itemize}

\subsection{Dockerfile -- Exemple pratique}

Voici un exemple de \texttt{Dockerfile} pour une application Python FastAPI :

\begin{lstlisting}[language=bash, caption=Dockerfile pour une API Python]
# Image de base legere
FROM python:3.11-slim

# Repertoire de travail
WORKDIR /app

# Installation des dependances
COPY requirements.txt .
RUN pip install --no-cache-dir -r requirements.txt

# Copie du code source
COPY . .

# Port expose
EXPOSE 8000

# Commande de lancement
CMD ["uvicorn", "app.main:app", "--host", "0.0.0.0", "--port", "8000"]
\end{lstlisting}

\subsection{Docker Compose}

Docker Compose permet de définir et d'orchestrer plusieurs conteneurs via un fichier \texttt{docker-compose.yml} :

\begin{lstlisting}[language=bash, caption=Exemple docker-compose.yml]
version: '3.9'
services:
  api:
    build: .
    ports:
      - "8000:8000"
    depends_on:
      - db
    environment:
      - DATABASE_URL=postgresql://user:pass@db:5432/mydb

  db:
    image: postgres:15
    volumes:
      - pgdata:/var/lib/postgresql/data
    environment:
      - POSTGRES_USER=user
      - POSTGRES_PASSWORD=pass
      - POSTGRES_DB=mydb

volumes:
  pgdata:
\end{lstlisting}

\subsection{Commandes Docker essentielles}

\begin{lstlisting}[language=bash, caption=Commandes Docker courantes]
# Construction d'une image
docker build -t monapp:v1 .

# Lancement d'un conteneur
docker run -d -p 8000:8000 --name monapp monapp:v1

# Lister les conteneurs actifs
docker ps

# Voir les logs
docker logs monapp

# Entrer dans un conteneur
docker exec -it monapp /bin/bash

# Pousser vers Docker Hub
docker push username/monapp:v1
\end{lstlisting}

% ==================== CLOUD COMPUTING ====================
\section{Cloud Computing}

\subsection{Définition}

Le cloud computing désigne la mise à disposition de ressources informatiques (calcul, stockage, réseau) via Internet, à la demande, avec facturation à l'usage.

\subsection{Modèles de service}

\begin{table}[H]
\centering
\begin{tabular}{@{}llp{6cm}@{}}
\toprule
\textbf{Modèle} & \textbf{Acronyme} & \textbf{Description} \\
\midrule
Infrastructure as a Service & IaaS & VMs, réseau, stockage brut (ex : AWS EC2, GCP Compute Engine) \\
Platform as a Service & PaaS & Environnement de développement géré (ex : Heroku, GCP App Engine) \\
Software as a Service & SaaS & Applications complètes en ligne (ex : Gmail, Salesforce) \\
Function as a Service & FaaS & Exécution serverless (ex : AWS Lambda, GCP Cloud Functions) \\
\bottomrule
\end{tabular}
\caption{Modèles de service Cloud}
\end{table}

\subsection{Modèles de déploiement}

\begin{itemize}
    \item \textbf{Cloud public} : Ressources partagées, gérées par un fournisseur (AWS, Azure, GCP)
    \item \textbf{Cloud privé} : Infrastructure dédiée à une organisation
    \item \textbf{Cloud hybride} : Combinaison public + privé
    \item \textbf{Multi-cloud} : Utilisation de plusieurs fournisseurs Cloud
\end{itemize}

\subsection{Principaux fournisseurs Cloud}

\begin{table}[H]
\centering
\begin{tabular}{@{}lll@{}}
\toprule
\textbf{Fournisseur} & \textbf{Produit Container} & \textbf{Orchestration} \\
\midrule
Amazon Web Services & Amazon ECS / EKS & Kubernetes managé \\
Google Cloud Platform & Google GKE & Kubernetes natif \\
Microsoft Azure & Azure AKS & Kubernetes managé \\
\bottomrule
\end{tabular}
\caption{Fournisseurs Cloud et services Kubernetes managés}
\end{table}

% ==================== KUBERNETES ====================
\section{Orchestration avec Kubernetes}

\subsection{Présentation}

Kubernetes (K8s) est un système open-source d'orchestration de conteneurs développé par Google. Il automatise le déploiement, la mise à l'échelle et la gestion des applications conteneurisées.

\subsection{Architecture Kubernetes}

\subsubsection{Plan de contrôle (Control Plane)}
\begin{itemize}
    \item \textbf{kube-apiserver} : Point d'entrée de l'API
    \item \textbf{etcd} : Base de données clé-valeur distribuée (état du cluster)
    \item \textbf{kube-scheduler} : Planification des Pods sur les nœuds
    \item \textbf{kube-controller-manager} : Gestion des contrôleurs (ReplicaSet, etc.)
\end{itemize}

\subsubsection{Nœuds de travail (Worker Nodes)}
\begin{itemize}
    \item \textbf{kubelet} : Agent sur chaque nœud, communique avec le Control Plane
    \item \textbf{kube-proxy} : Gestion des règles réseau
    \item \textbf{Container Runtime} : Docker / containerd
\end{itemize}

\subsection{Objets Kubernetes fondamentaux}

\begin{description}
    \item[Pod] Unité de base : un ou plusieurs conteneurs partageant réseau et stockage
    \item[Deployment] Gère le déploiement et la mise à jour des Pods (rolling update)
    \item[Service] Expose les Pods via une IP stable (ClusterIP, NodePort, LoadBalancer)
    \item[Ingress] Routage HTTP/HTTPS vers les Services
    \item[ConfigMap] Variables de configuration non-sensibles
    \item[Secret] Données sensibles chiffrées (mots de passe, tokens)
    \item[HPA] Horizontal Pod Autoscaler -- mise à l'échelle automatique
    \item[PersistentVolume] Stockage persistant découplé des Pods
\end{description}

\subsection{Manifestes YAML -- Exemples}

\begin{lstlisting}[language=bash, caption=Deployment Kubernetes]
apiVersion: apps/v1
kind: Deployment
metadata:
  name: monapp-deployment
  namespace: production
spec:
  replicas: 3
  selector:
    matchLabels:
      app: monapp
  template:
    metadata:
      labels:
        app: monapp
    spec:
      containers:
      - name: monapp
        image: username/monapp:v1
        ports:
        - containerPort: 8000
        resources:
          requests:
            cpu: "100m"
            memory: "128Mi"
          limits:
            cpu: "500m"
            memory: "512Mi"
\end{lstlisting}

\begin{lstlisting}[language=bash, caption=Service et HPA]
# Service
apiVersion: v1
kind: Service
metadata:
  name: monapp-service
spec:
  selector:
    app: monapp
  ports:
  - port: 80
    targetPort: 8000
  type: LoadBalancer

---
# Autoscaling
apiVersion: autoscaling/v2
kind: HorizontalPodAutoscaler
metadata:
  name: monapp-hpa
spec:
  scaleTargetRef:
    apiVersion: apps/v1
    kind: Deployment
    name: monapp-deployment
  minReplicas: 2
  maxReplicas: 10
  metrics:
  - type: Resource
    resource:
      name: cpu
      target:
        type: Utilization
        averageUtilization: 70
\end{lstlisting}

\subsection{Commandes kubectl essentielles}

\begin{lstlisting}[language=bash, caption=Commandes kubectl]
# Appliquer un manifest
kubectl apply -f deployment.yaml

# Voir l'etat des pods
kubectl get pods -n production

# Details d'un pod
kubectl describe pod <pod-name>

# Logs d'un pod
kubectl logs <pod-name> --follow

# Mise a l'echelle manuelle
kubectl scale deployment monapp-deployment --replicas=5

# Rollout d'une nouvelle version
kubectl set image deployment/monapp-deployment monapp=username/monapp:v2

# Rollback
kubectl rollout undo deployment/monapp-deployment
\end{lstlisting}

% ==================== CI/CD ====================
\section{CI/CD \& DevOps}

\subsection{Pipeline CI/CD}

Un pipeline CI/CD (\textit{Continuous Integration / Continuous Deployment}) automatise le cycle de vie d'une application : du code source jusqu'au déploiement en production.

\begin{keybox}{Étapes d'un pipeline CI/CD}
\textbf{1. Source} $\rightarrow$ \textbf{2. Build} $\rightarrow$ \textbf{3. Test} $\rightarrow$ \textbf{4. Package} $\rightarrow$ \textbf{5. Deploy}
\end{keybox}

\subsection{Exemple avec GitHub Actions}

\begin{lstlisting}[language=bash, caption=Pipeline GitHub Actions -- CI/CD vers GKE]
name: CI/CD Pipeline

on:
  push:
    branches: [main]

jobs:
  build-and-deploy:
    runs-on: ubuntu-latest
    steps:
      - uses: actions/checkout@v4

      - name: Build Docker image
        run: docker build -t gcr.io/$PROJECT_ID/monapp:$GITHUB_SHA .

      - name: Push to GCR
        run: docker push gcr.io/$PROJECT_ID/monapp:$GITHUB_SHA

      - name: Deploy to GKE
        run: |
          kubectl set image deployment/monapp \
            monapp=gcr.io/$PROJECT_ID/monapp:$GITHUB_SHA
\end{lstlisting}

% ==================== PROJET PRATIQUE ====================
\section{Application Pratique -- Smart Hardware Detector}

\subsection{Présentation du projet}

Dans le cadre de ce module, nous avons déployé le projet \textbf{Smart Hardware Detector} -- un classificateur d'images de matériel informatique -- sur une infrastructure Cloud conteneurisée.

\subsection{Architecture déployée}

\begin{keybox}{Stack technique}
\begin{itemize}[noitemsep]
    \item \textbf{Application} : FastAPI + PyTorch (ResNet18)
    \item \textbf{Conteneurisation} : Docker (image multi-stage)
    \item \textbf{Orchestration} : Kubernetes sur GKE (Google Kubernetes Engine)
    \item \textbf{CI/CD} : GitHub Actions
    \item \textbf{Monitoring} : Prometheus + Grafana
    \item \textbf{Scaling} : HPA (Horizontal Pod Autoscaler)
\end{itemize}
\end{keybox}

\subsection{Résultats du déploiement}

\begin{table}[H]
\centering
\begin{tabular}{@{}ll@{}}
\toprule
\textbf{Métrique} & \textbf{Valeur} \\
\midrule
Temps de démarrage conteneur & $< 3$ secondes \\
Replicas en production & 3 (min: 2, max: 10) \\
Disponibilité cible & 99.9\% \\
Stratégie de déploiement & Rolling Update \\
Image Docker (taille) & $\approx$ 850 Mo \\
\bottomrule
\end{tabular}
\caption{Métriques de déploiement}
\end{table}

% ==================== CONCLUSION ====================
\section{Conclusion}

Ce module nous a permis d'acquérir une vision complète de l'écosystème Cloud moderne :

\begin{itemize}
    \item La \textbf{virtualisation} constitue le socle technologique permettant l'isolation et la mutualisation des ressources.
    \item La \textbf{conteneurisation} avec Docker apporte portabilité et légèreté aux applications.
    \item \textbf{Kubernetes} orchestre les conteneurs à grande échelle avec haute disponibilité et auto-scaling.
    \item Le \textbf{Cloud Computing} offre une infrastructure élastique et une facturation à l'usage.
    \item Les pipelines \textbf{CI/CD} automatisent et fiabilisent les cycles de déploiement.
\end{itemize}

La maîtrise de ces technologies est aujourd'hui indispensable pour tout ingénieur souhaitant déployer des applications robustes, scalables et résilientes en production.

% ==================== BIBLIOGRAPHIE ====================
\begin{thebibliography}{9}

\bibitem{docker}
Docker Inc., \textit{Docker Documentation}, \url{https://docs.docker.com}, 2024.

\bibitem{kubernetes}
The Kubernetes Authors, \textit{Kubernetes Documentation}, \url{https://kubernetes.io/docs}, 2024.

\bibitem{gcp}
Google Cloud, \textit{Google Kubernetes Engine Documentation}, \url{https://cloud.google.com/kubernetes-engine/docs}, 2024.

\bibitem{tanenbaum}
A. Tanenbaum, \textit{Modern Operating Systems}, 4th ed., Pearson, 2014.

\bibitem{burns}
B. Burns et al., \textit{Kubernetes: Up and Running}, 3rd ed., O'Reilly Media, 2022.

\bibitem{mouat}
A. Mouat, \textit{Using Docker}, O'Reilly Media, 2016.

\end{thebibliography}

\end{document}
